\documentclass[aspectratio=169]{beamer}

\usetheme{Madrid}
\usecolortheme{default}

\usepackage{amsmath}
\usepackage{graphicx}
\usepackage{hyperref}

\title{ECE 251: Computer Architecture}
\subtitle{Week 13: Memory Hierarchies (Part 1)}
\author{Prof Rob Marano}
\date{Spring 2026}

% Cooper Union Brand Identity
\definecolor{CooperRed}{HTML}{8A1538}
\definecolor{CooperDark}{HTML}{5C0E25}
\setbeamercolor{structure}{fg=CooperRed}
\setbeamercolor{palette primary}{bg=CooperRed,fg=white}
\setbeamercolor{palette secondary}{bg=CooperRed,fg=white}
\setbeamercolor{palette tertiary}{bg=CooperDark,fg=white}
\setbeamercolor{title in head/foot}{bg=CooperDark,fg=white}
\logo{
\includegraphics[height=1cm]{../../images/cooper_union_logo.png}}
\titlegraphic{
\includegraphics[height=1.5cm,keepaspectratio]{../../images/cooper_union_logo.png}}

\begin{document}

\begin{frame}
  \titlepage
\end{frame}

\begin{frame}{Agenda}
  \tableofcontents
\end{frame}

\section{1. Introduction to Memory Hierarchy}

\begin{frame}{1. The Processor/Memory Paradox}
\begin{itemize}
    \item Computer architects face a fundamental physical paradox:
    \item \textbf{Processors are incredibly fast} (Clock cycles in nanoseconds or picoseconds).
    \item \textbf{Memory is relatively slow} (Access times in tens of nanoseconds or milliseconds).
    \item If the CPU has to wait for Main Memory on every instruction fetch and data load, the massive speed gains of pipelining ($CPI \approx 1.0$) are entirely lost!
\end{itemize}
\end{frame}

\begin{frame}{1. The Principle of Locality}
\begin{itemize}
    \item To keep the processor fed with data without stalling, we rely on the \textbf{Principle of Locality}:
    \vspace{0.5cm}
    \item \textbf{Temporal Locality (Locality in Time):}
    \begin{itemize}
        \item If an item is referenced, it will tend to be referenced again soon.
        \item \textit{Example:} Loops iterating over the same variables.
    \end{itemize}
    \vspace{0.3cm}
    \item \textbf{Spatial Locality (Locality in Space):}
    \begin{itemize}
        \item If an item is referenced, items whose addresses are close by will tend to be referenced soon.
        \item \textit{Example:} Sequentially accessing elements in an array.
    \end{itemize}
\end{itemize}
\end{frame}

\begin{frame}{1. The Memory Hierarchy Pyramid}
\begin{center}
    \includegraphics[height=5.5cm,keepaspectratio]{../../../../Image Bank/ch005-9780128201091/jpg-9780128201091/005001.jpg}
\end{center}
\begin{itemize}
    \item We exploit locality by implementing multiple tiers of memory.
    \item \textbf{Closer to CPU}: Faster, smaller, more expensive per bit.
    \item \textbf{Further from CPU}: Slower, larger, cheaper per bit.
\end{itemize}
\end{frame}

\begin{frame}{1. Key Terminology}
\begin{itemize}
    \item \textbf{Block (or Line):} The minimum unit of information that can be either present or not present in a two-level hierarchy.
    \item \textbf{Hit:} Data requested by the processor appears in some block in the upper level.
    \begin{itemize}
        \item \textbf{Hit Rate:} The fraction of memory accesses found in the upper level.
        \item \textbf{Hit Time:} The time required to access data in the upper level.
    \end{itemize}
    \item \textbf{Miss:} Data is not found in the upper level.
    \begin{itemize}
        \item \textbf{Miss Penalty:} The time required to fetch the block from the lower level into the upper level, plus the time to deliver it to the processor.
    \end{itemize}
\end{itemize}
\end{frame}

\section{2. Memory Technologies}

\begin{frame}{2. Memory Technologies: SRAM vs DRAM}
\begin{itemize}
    \item \textbf{SRAM (Static RAM):}
    \begin{itemize}
        \item Used for caches (L1, L2, L3).
        \item Fast (0.5ns - 2.5ns access time).
        \item Requires no refreshing to keep data (hence ``static'').
        \item Low density and very expensive (6 transistors per bit).
    \end{itemize}
    \vspace{0.3cm}
    \item \textbf{DRAM (Dynamic RAM):}
    \begin{itemize}
        \item Used for Main Memory.
        \item Slower (50ns - 70ns access time).
        \item Data stored as a charge on a capacitor; must be periodically ``refreshed'' because charge leaks.
        \item High density and cheaper (1 transistor per bit).
    \end{itemize}
\end{itemize}
\end{frame}

\begin{frame}{2. Non-Volatile Memory}
\begin{itemize}
    \item \textbf{Flash Memory:}
    \begin{itemize}
        \item Non-volatile semiconductor memory (EEPROM).
        \item Survives power-off.
        \item Reads are fast, but writes are significantly slower and physically wear out the memory cells over time (requires wear leveling).
    \end{itemize}
    \vspace{0.3cm}
    \item \textbf{Magnetic Disk:}
    \begin{itemize}
        \item Used for massive secondary storage.
        \item Relies on moving mechanical parts (platters and read/write heads).
        \item Measured in milliseconds (millions of times slower than SRAM).
    \end{itemize}
\end{itemize}
\end{frame}

\section{3. The Basics of Caches}

\begin{frame}{3. The Cache}
\begin{itemize}
    \item The \textbf{Cache} is the level of the memory hierarchy closest to the CPU.
    \item The simplest way to map memory blocks into the cache is a \textbf{Direct-Mapped Cache}.
    \item \textbf{Direct-Mapped Caching:} Each memory block maps to exactly \textit{one} specific block in the cache.
    \item \textbf{Mapping Formula:} \texttt{(Block address) modulo (Number of blocks in the cache)}
    \item Since the number of cache blocks is usually a power of 2, the modulo operation is just the lower $n$ bits of the block address.
\end{itemize}
\end{frame}

\begin{frame}{3. Direct Mapped Cache Architecture}
\begin{center}
    \includegraphics[height=6.5cm,keepaspectratio]{../../../../Image Bank/ch005-9780128201091/jpg-9780128201091/005007.jpg}
\end{center}
\end{frame}

\begin{frame}{3. Tags and Valid Bits}
\begin{itemize}
    \item Because multiple memory blocks map to the same cache index, we need a way to know exactly \textit{which} memory block is currently sitting in the cache.
    \vspace{0.3cm}
    \item \textbf{Tag:} The upper portion of the address is stored alongside the data. The cache compares the requested address's tag against the stored tag to verify a match.
    \vspace{0.3cm}
    \item \textbf{Valid Bit:} A single bit added to the cache block to indicate whether the block contains valid data (1) or garbage/empty space (0).
\end{itemize}
\end{frame}

\begin{frame}{3. Handling Cache Misses}
\begin{itemize}
    \item When a miss occurs, the control unit must:
    \begin{enumerate}
        \item \textbf{Stall} the CPU pipeline (suspend execution).
        \item \textbf{Fetch} the requested block from the next lower level of the memory hierarchy (Main Memory).
        \item \textbf{Write} the block into the cache (updating the tag and setting the valid bit).
        \item \textbf{Restart} the instruction that caused the miss.
    \end{enumerate}
\end{itemize}
\end{frame}

\begin{frame}{3. Block Size Considerations}
\begin{itemize}
    \item Increasing the block size takes advantage of spatial locality, reducing the miss rate.
    \item \textbf{The Trade-off:}
    \begin{itemize}
        \item If the block size gets too large compared to the overall cache size, the miss rate can actually \textit{increase} because fewer total blocks can be held in the cache simultaneously (increasing conflict/capacity misses).
        \item A larger block size increases the \textit{miss penalty} since more data must be fetched from slower memory over the system bus.
    \end{itemize}
\end{itemize}
\end{frame}

\begin{frame}
\begin{center}
    \Huge Questions?
\end{center}
\end{frame}

\end{document}
